\section{Methodology}

We describe our benchmarking pipeline for image classification on Caltech-101. Given a dataset of images $\{\mathbf{I}_i\}_{i=1}^{N}$ with class labels $\{y_i\}_{i=1}^{N}$, our objective is to train and evaluate multiple classification models under consistent settings. The pipeline comprises four components: (1) stratified data preparation, (2) model architectures with ImageNet pretraining, (3) fine-tuning and training protocols, and (4) evaluation metrics. All implementation details are aligned with our released codebase.

\subsection{Dataset Preparation}

Caltech-101~\citep{fei2004learning} contains 9{,}144 images across 102 categories (101 object classes plus a background class), with notable class imbalance ranging from approximately 40 to 800 images per class. We include the background category in our experiments (i.e., \texttt{exclude\_background = false}).

We apply a stratified split to preserve class proportions in each subset. For each class $c$ with $n_c$ images, we partition into train, validation, and test sets:
\vspace{3em}
\begin{align}
\eqnmarkbox[Emerald]{ntrain}{n_c^{\text{train}}} &= \max\!\Big(1,\;\left\lfloor \eqnmarkbox[OliveGreen]{rtrain}{0.70} \cdot \eqnmarkbox[NavyBlue]{nc}{n_c} \right\rfloor\Big), \quad
\eqnmarkbox[WildStrawberry]{nval}{n_c^{\text{val}}} = \left\lfloor \eqnmarkbox[BurntOrange]{rval}{0.15} \cdot n_c \right\rfloor, \nonumber\\
\eqnmarkbox[Plum]{ntest}{n_c^{\text{test}}} &= n_c - n_c^{\text{train}} - n_c^{\text{val}},
\label{eq:split}
\end{align}
\annotate[yshift=1.8em]{above,left}{ntrain}{\tiny Train count}
\annotate[yshift=1.8em]{above,left}{rtrain}{\tiny Train ratio}
\annotate[yshift=1.8em]{above,right}{nc}{\tiny Class size}
\annotate[yshift=1.8em]{above,right}{nval}{\tiny Val count}
\annotate[yshift=1.8em]{above,right}{rval}{\tiny Val ratio}
\annotate[yshift=-1em]{below,right}{ntest}{\tiny Test (remainder)}
\vspace{1.5em}

with a minimum of one training sample per class enforced via the $\max(1, \cdot)$ operator. We use a fixed random seed ($s = 42$) for reproducibility. The resulting split yields 6{,}352 / 1{,}326 / 1{,}466 samples for train / validation / test, with class proportions preserved across all subsets.

\subsection{Model Architectures}

We benchmark five methods spanning three paradigms: convolutional networks, vision transformers, and classical machine learning. All deep models use ImageNet-pretrained weights; only the final classification head is replaced with a new linear layer mapping to $C = 102$ classes.

\begin{itemize}[nosep,leftmargin=*]
\item \textbf{ResNet-50}~\citep{he2016deep} learns residual mappings $\mathcal{F}(\mathbf{x}) + \mathbf{x}$ via skip connections. We load \texttt{IMAGENET1K\_V2} weights and replace the final \texttt{fc} layer ($2048 \to C$).

\item \textbf{EfficientNet-B2}~\citep{tan2019efficientnet} uses compound scaling to jointly optimize depth, width, and resolution. We load \texttt{IMAGENET1K\_V1} weights and replace \texttt{classifier[1]} ($1408 \to C$).

\item \textbf{ViT-B/16}~\citep{dosovitskiy2020image} splits each image into non-overlapping $16 \times 16$ patches and processes them through 12 transformer layers. We load \texttt{IMAGENET1K\_V1} weights and replace \texttt{heads.head} ($768 \to C$).

\item \textbf{ConvNeXt-Tiny}~\citep{liu2022convnet} modernizes standard ConvNets with design elements from transformers (larger kernels, layer normalization, inverted bottlenecks). We load \texttt{IMAGENET1K\_V1} weights and replace \texttt{classifier[2]} ($768 \to C$).

\item \textbf{SVM + ResNet-18 features}~\citep{donahue2014decaf} serves as a classical baseline. We extract 512-dimensional feature vectors from a frozen, ImageNet-pretrained ResNet-18 by removing its final FC layer and global-average pooling the output. Features are standardized with zero mean and unit variance, then fed to an RBF-kernel SVM ($C = 1$, $\gamma = \text{scale}$). No end-to-end fine-tuning is performed.
\end{itemize}

\subsection{Training Details}

\textbf{Differential learning rates.} For each deep model, we partition parameters into backbone (pretrained) and head (randomly initialized) groups with separate learning rates. For ResNet-50, EfficientNet-B2, and ConvNeXt-Tiny we use $\eta_{\text{bb}} = 10^{-4}$ and $\eta_{\text{head}} = 10^{-3}$. For ViT-B/16, which is more sensitive to learning rate, we use $\eta_{\text{bb}} = 2 \times 10^{-5}$ and $\eta_{\text{head}} = 2 \times 10^{-4}$. All models are optimized with AdamW~\citep{loshchilov2019decoupled} ($\lambda = 10^{-4}$).

\textbf{Learning rate schedule.} We apply a linear warmup over the first $T_w = 5$ epochs (ramping from $0.01 \eta$ to $\eta$), followed by cosine annealing:
\begin{equation}
\eta_t = \eta_{\min} + \tfrac{1}{2}(\eta_{\max} - \eta_{\min})\!\left(1 + \cos \tfrac{(t - T_w)\pi}{T - T_w}\right),
\end{equation}
where $T = 30$ is the maximum number of epochs. Training is terminated early if validation accuracy does not improve for 5 consecutive epochs.

\textbf{Regularization.} We apply label smoothing~\citep{szegedy2016rethinking} ($\epsilon = 0.1$) in the cross-entropy loss and gradient clipping (max norm 1.0). Mixed-precision training (FP16) is enabled on GPU.

\textbf{Data augmentation.} During training: resize to $256 \times 256$, random crop to $224 \times 224$, random horizontal flip, color jitter (brightness $= 0.4$, contrast $= 0.4$, saturation $= 0.4$, hue $= 0.1$), and random erasing~\citep{zhong2020random} ($p = 0.25$). At evaluation: resize to $256 \times 256$ and center crop to $224 \times 224$. All images are normalized with ImageNet statistics ($\mu = [0.485, 0.456, 0.406]$, $\sigma = [0.229, 0.224, 0.225]$).

\textbf{Loss function.} We minimize the cross-entropy loss with label smoothing:
\vspace{2em}
\begin{equation}
\label{eq:loss}
\mathcal{L} = -\sum_{k=1}^{\eqnmarkbox[NavyBlue]{K}{K}} \left[ \left(1-\eqnmarkbox[OliveGreen]{eps}{\epsilon}\right)\, \mathbb{1}[y=k] + \frac{\epsilon}{\eqnmarkbox[NavyBlue]{K}{K}} \right] \log \eqnmarkbox[WildStrawberry]{pk}{p_k},
\end{equation}
\annotate[yshift=1.8em]{above,left}{K}{\tiny Num.\ classes}
\annotate[yshift=1.8em]{above,right}{eps}{\tiny Label smoothing}
\annotate[yshift=-1em]{below,right}{pk}{\tiny Predicted prob.}
\vspace{1em}

\noindent where $K = 102$, $\epsilon = 0.1$, and $p_k$ is the softmax probability for class $k$. We use batch size 32 with 4 data-loading workers and pinned memory.

\subsection{Evaluation Protocol}

All metrics are computed on the held-out test set. We report: (1)~overall accuracy, (2)~top-5 accuracy (for deep models), (3)~macro- and weighted-averaged precision, recall, and F1-score, (4)~per-class accuracy (i.e., per-class recall), and (5)~confusion matrices. Macro averaging treats all classes equally regardless of sample count, while weighted averaging accounts for class imbalance. Per-class accuracy reveals categories where models underperform; confusion matrices identify the most common inter-class confusions.
